\section{The algebraic model and the decision threshold}
\label{sec:algebraic-model}

The input is an explicitly listed, validated classical knot diagram with
$n$ crossings. Complexity statements use normalized edge labels of
$O(\log(n+2))$ bits; reading and normalizing a longer external encoding
incurs its actual input-size cost. The empty input denotes one
crossing-free circle. A scan
order determines, after each crossing, a finite set of live boundary
labels. Its cardinality is even. As in the preceding implementation, the
objects of the scanning category are loopless matchings of these labels;
closed circles are replaced by their two summands by delooping. We retain
homological degrees and work over $\F$, but forget the quantum grading.
In particular, the basis changes below are assertions in the
quantum-ungraded category. A quantum-graded extension would require
additional compatibility conditions on those basis changes.

Write
\[
 r_2(K)=\dim_{\F}\Kh(K;\F).
\]
The decision criterion used by the existing recognizer is
\begin{equation}\label{eq:rank-two}
 K\text{ is an unknot}\quad\Longleftrightarrow\quad r_2(K)=2.
\end{equation}
Here is the coefficient justification. The reduced/unreduced splitting
over $\F$ gives $r_2(K)=2\dim_{\F}\rKh(K;\F)$
\cite{shumakovitch}. The universal coefficient theorem implies
\[
 \dim_{\Q}\rKh(K;\Q)\le\dim_{\F}\rKh(K;\F).
\]
The rational-rank unknot-detection theorem then excludes every nontrivial
knot when the reduced $\F$ rank is one~\cite{km}. The reverse implication
follows from the unknot complex. This argument does not assume that a
rank-one reduction modulo two excludes odd integral torsion.

We use the notation
\[
 \capthree(a)=\min(3,a),\qquad a\in\N.
\]
A completed capped value of two therefore certifies \UNKNOT, while a
value of three certifies \KNOTTED. For a valid knot the latter value means
that the actual rank is at least four. Values zero and one indicate a
violated input promise or an implementation invariant; they are not
interpreted as topological verdicts.

The algebraic correctness statements are relative to the established
Bar--Natan gluing, delooping, and cancellation operations
\cite{khovanov,barnatan}. The new algorithms operate on their finite
matrices and use exact binary arithmetic. They do not constitute a formal
verification of the entire topological foundation.

\section{Exact interval normalization of square-zero components}
\label{sec:interval-normalization}

\subsection{A recognizable coefficient subcategory}

For a matching $m$ with $k$ arcs, the endomorphism ring used by the scanner
is the dot algebra
\begin{equation}\label{eq:dot-algebra}
 R_m=\End(m)\cong
 \F[x_1,\ldots,x_k]/(x_1^2,\ldots,x_k^2).
\end{equation}
Its binary basis consists of the squarefree dot monomials; its dimension
is $2^k$. The identity is the constant monomial. This is the familiar
diagonal endomorphism algebra of the matching category
\cite{khovanov,barnatan}.

Let $\mathfrak m=(x_1,\ldots,x_k)$ be the augmentation ideal. Every
$\theta\in\mathfrak m$ satisfies $\theta^2=0$: in characteristic two,
the square of a sum is the sum of its squares, and each nonconstant
squarefree monomial has square zero. Conversely, an element with constant
coefficient one is a unit. Thus the nonzero augmentation elements are
exactly the nonzero nonunits of this ring.

The undirected differential graph has one vertex per object and an edge
when a differential entry is nonzero. Its connected components are
literal direct summands of the complex. The earlier sharing backend
identifies repeated components by complete serialization~\cite{prior}.
The following condition exposes additional summands inside such a
component.

\begin{definition}[Pure component]\label{def:pure}
A nonsingleton differential component is \emph{pure} if all its objects
have one matching $m$ and every nonzero differential entry is one fixed
nonzero $\theta\in R_m$ satisfying $\theta^2=0$. A singleton is pure
without a choice of $\theta$.
\end{definition}

Purity is a test of the actual stored matrix. We do not posit an oracle
which determines whether an arbitrary complex can be made pure. In a pure
component, write the degree-$h$ term as $m^{d_h}$. Then
\begin{equation}\label{eq:theta-matrix}
 \delta_h=\theta M_h,\qquad
 M_h:\F^{d_h}\longrightarrow\F^{d_{h+1}},
\end{equation}
where $\delta_h$ is the differential and $M_h$ is a binary matrix.
The matrices $M_h$ need not satisfy
$M_{h+1}M_h=0$. The actual differential squares to zero because
$\theta^2=0$. The distinction is essential for the algorithm: its rank
calculations use products of the $M_h$, not products of the actual
cobordism differentials.

\begin{definition}[Interval complex]
For integers $a\le b$, let $I_{[a,b]}(m,\theta)$ have one copy of $m$ in
each degree $a,\ldots,b$, with differential $\theta$ between successive
copies and zero elsewhere. Its length is $b-a+1$. After removing the birth
shift $a$, write this complex as $I_{b-a+1}(m,\theta)$.
\end{definition}

\subsection{The exact decomposition and its computation}

\begin{theorem}[Exact interval normalization]\label{thm:interval}
Every pure component is chain-isomorphic to a direct sum of interval
complexes. The isomorphism is given by binary changes of basis within
each homological degree. The interval multiplicities are uniquely
determined by the binary matrices in~\eqref{eq:theta-matrix}.
\end{theorem}

\begin{proof}
We first prove the elementary interval decomposition for an arbitrary
finite sequence of binary vector spaces and maps $M_h$; this is the
equioriented type-$A$ case of the usual interval classification
\cite{carlsson}. Suppose the sequence is nonzero. Let $a$ be its earliest
nonzero degree, and choose the largest $b\ge a$ for which the composite
from $V_a$ to $V_b$ is nonzero, counting the identity when $b=a$.
Choose $v_a$ whose image $v_b$ is nonzero, and let $v_h$ be its successive
images for $a\le h\le b$. Choose a covector $\lambda_b$ with
$\lambda_b(v_b)=1$, and define
\[
 \lambda_h=\lambda_bM_{b-1}\cdots M_h,
 \qquad p_h=v_h\lambda_h\quad(a\le h\le b).
\]
Outside this interval put $p_h=0$. We have $\lambda_h(v_h)=1$, so each
$p_h$ is idempotent. For an interior arrow,
$p_{h+1}M_h=M_hp_h$. At the left endpoint $V_{a-1}=0$; at the right
endpoint $M_bv_b=0$ by maximality of $b$. Hence $p$ is an idempotent
endomorphism of the entire sequence. Its image is one interval and its
kernel is an invariant complement. Induction on the total vector-space
dimension proves existence.

Every binary invertible change-of-basis matrix remains invertible after
its entries are multiplied by $\mathrm{id}_m$ in the additive category:
its binary inverse supplies the actual inverse. Applying the changes
of basis to $\theta M_h$ therefore gives a chain isomorphism of the
original component with the corresponding $\theta$-intervals.
The multiplicity formula below proves uniqueness.
\end{proof}

Let the occupied degree range be $u,\ldots,v$. For $u\le a\le b\le v$
define the rank invariant
\begin{equation}\label{eq:rank-invariant}
 \rho(a,b)=\rank(M_{b-1}\cdots M_a),
 \qquad \rho(a,a)=d_a.
\end{equation}
Ranks with an endpoint outside the occupied range are assigned the value
zero. An interval $[p,q]$ contributes one to $\rho(a,b)$ exactly when
$p\le a\le b\le q$. Taking successive differences gives
\begin{equation}\label{eq:mobius-bars}
 \mu(a,b)=\rho(a,b)-\rho(a-1,b)
          -\rho(a,b+1)+\rho(a-1,b+1).
\end{equation}
This is the multiplicity of $[a,b]$. In particular, the result is a
nonnegative integer, and
\begin{equation}\label{eq:barcode-audit}
 \sum_{a\le b}\mu(a,b)(b-a+1)=\sum_h d_h.
\end{equation}
These identities give both a deterministic algorithm and useful
independent consistency checks. The implementation can compute all ranks
by binary matrix multiplication and elimination, form the differences,
and reconstruct the rank invariant from the resulting intervals.

If there are $H$ degree layers and $B$ objects, a conservative elementary
bound is $O(H^2B^3)$ binary field operations. No efficiency claim depends
on this exponent being optimal. Bit-packed matrices provide the practical
implementation. Recognition of purity and the required coefficient
comparisons are polynomial in the stored matrix and coefficient lengths.

Replacing the component by explicit intervals retains exactly $B$ objects
by~\eqref{eq:barcode-audit}. Its advantage is that these objects now lie
in separately shareable summands. Number an interval's vertices in
increasing degree, and remove its birth shift. Its complete template is
then determined by $(m,\theta,\text{length})$. Length-one templates do not
use $\theta$. Equal templates can be continued once, with exact shifted
multiplicities or with the existing saturated multiplicities.

\begin{example}[A strict improvement over graph-component sharing]
\label{ex:dense-pure}
Take $B\ge2$ copies of $m$ in degree zero and $B$ copies in degree one,
and put $\theta$ in every differential entry. Its differential graph is
the connected graph $K_{B,B}$, and no entry is a unit. Literal component
sharing leaves one template with $2B$ objects. The scalar matrix is the
all-ones matrix of rank one. Interval normalization produces one
length-two interval and $B-1$ singletons in each of the two degrees.
After shifted sharing, only three objects remain: two in the interval
and one singleton with multiplicity polynomial $(B-1)(1+z)$.

This is a valid algebraic stress family. It is not asserted that arbitrary
members occur as intermediate states of the current knot scanner. That
additional realizability assertion would need a separate construction.
\end{example}

The replacement must apply to an entire direct summand. A tower attached
to other matchings may be normalized only if its basis changes are also
applied to every external incident differential. Simply discarding those
attachments would invalidate the complex. The exact splitter developed
in the next part of the article addresses this issue by producing and
checking an actual change of basis.

\section{Short intervals with identical capped continuation ranks}
\label{sec:continuation-quotient}

Exact normalization still leaves intervals of unbounded length. They
need not split further: finite free strings are indecomposable over the
dual numbers~\cite{keller}. The decision problem permits a different operation. A long
string can be replaced by a short one when every possible remaining
continuation has the same capped homology rank. Such a replacement is
not asserted to be a chain-homotopy equivalence.

\subsection{The dual-number tensor calculation}

Put $A=\F[\epsilon]/(\epsilon^2)$, and let $J_l$ be the complex of $l$
copies of the free rank-one $A$-module, in consecutive degrees, with every
differential multiplication by $\epsilon$. All tensor products of
complexes below are totalized, and all dimensions refer to total
$\F$-homology. Homological shifts do not affect these dimensions.

\begin{lemma}[Tensor pairing of intervals]\label{lem:interval-tensor}
For positive integers $l,m$,
\[
 \dim_{\F}H(J_l\otimes_A J_m)=2\min(l,m),
 \qquad
 \dim_{\F}H(J_l\otimes_A\F)=l.
\]
\end{lemma}

\begin{proof}
For the first formula, identify the $lm$ free basis positions of the
total tensor product, after forgetting their grading, with
\[
 R=\F[u,v]/(u^l,v^m).
\]
The total differential on $A\otimes_{\F}R$ is multiplication by
$\epsilon$ times the linear operator $T$ given by multiplication by
$u+v$ on $R$. There are no sign distinctions in characteristic two.
This differential has binary rank $\rank(T)$, so its total homology has
dimension $2lm-2\rank(T)=2\dim\coker(T)$. But
\[
 \coker(T)\cong
 \F[u,v]/(u^l,v^m,u+v)
 \cong\F[u]/(u^{\min(l,m)}),
\]
which has dimension $\min(l,m)$. Collecting the degreewise matrices
into one graded total differential preserves the sum of their ranks.
For the second formula, $\epsilon$ acts by zero on the simple module,
so all differentials vanish and there are $l$ one-dimensional terms.
\end{proof}

The needed derived-category classification is particularly small. Every
bounded complex of finite $A$-modules is isomorphic in the bounded
derived category to a finite direct sum of shifts of the $J_m$ and shifts
of the simple module $\F$~\cite[Example 3.7(b)]{keller}.
Appendix~\ref{app:dual-number-classification} gives an elementary proof
using minimal free resolutions and the finite interval argument above.
We use this established classification as an algebraic input; no novelty
claim is made for it.

\begin{theorem}[The continuation rank formula]\label{thm:rank-formula}
Let $D$ be a bounded complex of finite-dimensional $A$-modules. There are
nonnegative integers $a,b_1,b_2,\ldots$, only finitely many nonzero, such
that, for every $l\ge1$,
\begin{equation}\label{eq:continuation-rank-formula}
 r_l(D):=\dim_{\F}H(J_l\otimes_A D)
   =a l+2\sum_{m\ge1}b_m\min(l,m).
\end{equation}
The integer $a$ counts the simple-module summands of $D$, and $b_m$
counts its free interval summands of length $m$, ignoring shifts.
\end{theorem}

\begin{proof}
The complex $J_l$ is bounded and free, hence tensoring with it preserves
quasi-isomorphisms and computes the derived tensor product. Decompose
$D$ as above and apply Lemma~\ref{lem:interval-tensor} to each summand.
The decomposition is finite, and homology dimensions add under direct
sums. This proves the formula simultaneously for all $l$.
\end{proof}

The theorem does not require the algorithm to compute the decomposition
of the potentially enormous suffix $D$. It proves an identity for every
such suffix, which justifies a local replacement before the suffix is
processed.

\begin{corollary}[Universal length cutoff]\label{cor:universal-cutoff}
For any positive integer threshold $c$, any $D$ as above, and any $l\ge1$,
\[
 \min(c,r_l(D))=\min(c,r_{\min(l,c)}(D)).
\]
In particular, length three is a safe cutoff for rank capped at three
under arbitrary bounded dual-number continuations.
\end{corollary}

\begin{proof}
Only $l\ge c$ needs consideration. If $a>0$, both ranks are at least
$c$. If $a=0$ and some $b_m>0$ has $m\ge c$, that summand contributes
at least $2c$ to both ranks. Otherwise every nonzero $b_m$ has $m<c$,
so the right side of~\eqref{eq:continuation-rank-formula} has stabilized
by $l=c$ and the two ranks are exactly equal.
\end{proof}

The cutoff in this fully general statement cannot be lowered uniformly:
for $D=\F$, the rank is $r_l(D)=l$. The geometric continuations of the
scanner have an extra parity property, which permits a sharper cutoff.

\subsection{Why a tangle continuation has the required module structure}

Fix a matching $m$ and a square-zero endomorphism $\theta$. Let $F$ mean
attach all remaining crossing complexes, close the diagram, and apply
the ordinary unreduced Frobenius TQFT. For the proof, use the strict raw
complex before optional Gaussian contractions. Planar gluing of complexes
is additive, preserves composition, and uses the total tensor
differential~\cite{barnatan}. Thus
\[
 D=F(m),\qquad \epsilon\cdot v=F(\theta)(v)
\]
defines a bounded complex of finite $A$-modules: $F(\theta)$ is a
degree-zero chain map and $F(\theta)^2=F(\theta^2)=0$. Inspecting the
totalized differential gives
\begin{equation}\label{eq:functor-tensor}
 F\bigl(I_l(m,\theta)\bigr)\cong J_l\otimes_A D.
\end{equation}
This identifies the balanced module action required by
Theorem~\ref{thm:rank-formula}. It is stronger than merely saying that
the continuation is a linear operation. The implementation may perform
homotopy-preserving contractions, but the raw functor supplies this
strict model for the proof.

\begin{lemma}[Even continuation rank]\label{lem:even-suffix}
If $m$ admits a nonzero square-zero endomorphism, then
$\dim_{\F}H(F(m))$ is even for every actual scanner continuation.
\end{lemma}

\begin{proof}
The empty matching has endomorphism ring $\F$, so $m$ has at least one
arc. Every complete smoothing of the suffix, glued to $m$, is a nonempty
union of circles. Its unreduced TQFT vector space has dimension $2^c$
for some $c\ge1$. Hence the raw total chain dimension of $F(m)$ is even.
For any finite complex of binary vector spaces,
\[
 \dim H(D)=\dim D-2\rank(d_D),
\]
so its total homology dimension has the same parity. This proves the
claim. In particular, the argument does not require a marked-point
implementation or an assumption that the $A$-module complex is free.
\end{proof}

\begin{theorem}[Geometric length cutoff two]\label{thm:geometric-cutoff}
For every matching $m$, nonzero square-zero $\theta\in\End(m)$, positive
length $l$, and actual remaining continuation $F$,
\begin{equation}\label{eq:length-two}
 \capthree\!\left(\dim H(F(I_l(m,\theta)))\right)
 =
 \capthree\!\left(\dim H(F(I_{\min(l,2)}(m,\theta)))\right).
\end{equation}
\end{theorem}

\begin{proof}
By Lemma~\ref{lem:even-suffix} and
\eqref{eq:continuation-rank-formula},
$r_1(D)=a+2\sum_m b_m$ is even, so $a$ is even. Suppose $l\ge2$.
If $a>0$, then $a\ge2$ and its contribution at length two is already
at least four. If $a=0$ but some $b_m>0$ has $m\ge2$, that summand
also contributes at least four at length two. In the remaining case,
only $b_1$ can be nonzero, and $r_l(D)=2b_1$ is independent of $l$.
These three cases prove the equality. Length one is unchanged.
\end{proof}

\begin{example}[The geometric cutoff is sharp]\label{ex:cutoff-sharp}
Let $m$ pair four cyclically ordered endpoints as
$(1,2)$ and $(3,4)$, and close it with the other crossingless pairing
$(1,4)$ and $(2,3)$. The closure is one circle. Write $x_1,x_2$ for dots
on the two arcs of $m$ and set $\theta=x_1+x_2$. Under closure both dots
act by the same circle operator $x$, so $F(\theta)=x+x=0$. The image
of $m$ is the two-dimensional circle space. Consequently
\[
 \dim H(F(I_l(m,\theta)))=2l.
\]
Replacing a length-two interval by a singleton would change a capped
value of three to two. Thus cutoff one cannot preserve these geometric
observations uniformly. This example concerns valid algebraic objects
and a crossingless planar closure; it does not claim that each source
interval is itself the full complex of a knot diagram.
\end{example}

\subsection{Compositional correctness of the decision backend}

\begin{definition}[Continuation equivalence]\label{def:continuation-equivalence}
For two complexes $C,C'$ on the same boundary, write $C\equiv_3 C'$ if
every permitted remaining crossing sequence and ordinary closure $F$
satisfies
\[
 \capthree\!\left(\dim H(F(C))\right)
 =\capthree\!\left(\dim H(F(C'))\right).
\]
\end{definition}

This is an equivalence relation by equality of the specified observations.
It is preserved by adjoining a common direct summand, taking finitely
many copies, and performing a common crossing extension: in the last
case, any further completion is also a permitted completion before that
extension. Theorem~\ref{thm:geometric-cutoff} identifies long intervals
with length-two representatives for this relation. It does not identify
their homotopy types.

The prior backend stores a template together with nonnegative
multiplicities, optionally recording absolute shifts in a polynomial
$p(z)\in\N[z]$~\cite{prior}. For total-rank decisions, replace the
weight by $\capthree(p(1))$. Addition and multiplication use saturation:
\[
 a\oplus b=\min(3,a+b),\qquad a\otimes b=\min(3,ab).
\]
The map $\N\to\{0,1,2,3\}$ given by saturation preserves both
operations. These are nonnegative counts, not coefficients reduced
modulo two.

\begin{theorem}[Decision correctness under repeated replacements]
\label{thm:decision-correctness}
Start with the existing exact scanning complex. At any stages, perform
exact unit cancellations, exact changes of basis, direct-summand
decompositions, and sharing of equal normalized templates. For each
pure summand, apply Theorem~\ref{thm:interval}; in decision mode shorten
each resulting interval to length $\min(l,2)$ and saturate its
multiplicity at three. Assume that structural choices do not depend on
the discarded multiplicities or absolute shifts. At final closure the
reported count is precisely $\capthree(r_2(K))$.
\end{theorem}

\begin{proof}
Fix the full future crossing order at any proposed replacement. By
Theorem~\ref{thm:geometric-cutoff}, completing a single replaced interval
would give the same capped total homology rank. Completion is additive
on independent summands and commutes with homological shifts. Since
capping commutes with addition and multiplication of nonnegative
integers, replacing any number of independent copies preserves the
capped rank of the completed whole complex. Exact changes of basis,
contractions, and component sharing preserve the actual homotopy type
before capping and hence preserve this observation as well.

Repeat this argument at every later replacement, comparing the
counterfactual full completion immediately before and immediately after
it. Transitivity shows that all replacements together preserve the
capped final rank. At closure, the empty matching has endomorphism
ring $\F$; complete elimination leaves isolated generators, whose
capped multiplicities compute that rank. Finally apply
\eqref{eq:rank-two} for the verdict.
\end{proof}

This proof deliberately uses an observational invariant after shortening.
The represented complex need not retain the original homotopy type,
exact rank, or degree-by-degree homology. The exact-rank backend therefore
uses interval splitting without shortening. The decision backend reports
a capped rank and does not supply exact homological counts or cycle
representatives. Likewise, exact signed Euler data from the former
representation cannot silently be carried through an interval shortening.

The current length alone does not give a completed rank bound greater
than two. For a geometric example, take one arc $m$ with $\theta=x$ and
close it to one circle. Then $D=A$ with $\epsilon$ acting by $x$, and
Lemma~\ref{lem:interval-tensor} gives
$\dim H(J_l\otimes_A A)=2$ for every $l$. Thus arbitrarily long legal
interval complexes can have a rank-two closure observation. For
arbitrary module continuations, an acyclic $D$ can even give rank zero.
The theorem says that the short representative has the same completed
threshold observation; it does not license a verdict from the current
number of stored objects.

\section{A checkpoint bound independent of component size}
\label{sec:checkpoint-complexity}

\subsection{The state count after normalization}

Fix the implemented ordering and exact splitting policy. Let $w$ be the
largest frontier size over the entire planned scan. Let a \emph{pure
checkpoint} be a stage at which, after unit elimination and the chosen
exact splitting operations, every differential component satisfies
Definition~\ref{def:pure}. Include the initial and final states among
the checkpoints. Let $g$ bound the number of crossings between
successive checkpoints. A checkpoint is recognizable from its current
matrix. A useful uniform bound on their gaps is an additional theorem
about an input family, not something implied by their definition.

For even $w\ge0$ put
\begin{equation}\label{eq:M-N}
 M(0)=1,\qquad M(w)=(w-1)!!\quad(w>0),
 \qquad N(w)=2^{\,2^{w/2}-1}-1.
\end{equation}
There are at most $M(w)$ frontier matchings. The quantity $N(w)$ counts
the nonzero augmentation elements of the endomorphism ring of a
matching with $w/2$ arcs: its $2^{w/2}-1$ nonconstant basis monomials
have arbitrary binary coefficients. In particular, $N(0)=0$.
Counting all perfect matchings avoids assuming that the frontier bounds
one disc or that numeric labels specify a useful cyclic order.

\begin{samepage}
\begin{proposition}[Physical size at a pure checkpoint]
\label{prop:checkpoint-size}
After interval normalization, length cutoff two, and complete sharing,
a pure checkpoint satisfies
\begin{equation}\label{eq:checkpoint-P}
\begin{aligned}
 \text{stored templates}&\le M(w)(1+N(w)),\\
 \text{stored objects}&\le P(w):=M(w)(1+2N(w)),
\end{aligned}
\end{equation}
independently of multiplicities and of the sizes of the
original connected differential components.
\end{proposition}
\end{samepage}

\begin{proof}
For each matching there is one possible singleton type, with one object.
Every other retained interval has length two and is determined by its
matching and one of the at most $N(w)$ possible nonzero square-zero
endomorphisms. It contributes two objects. Constructing the vertices
in degree order makes equal triples have equal complete serializations;
no graph-isomorphism oracle is needed for these types. Summing these
counts gives the bounds. A smaller actual frontier only decreases
both $M$ and $N$.
\end{proof}

For comparison, the exact-rank version retains lengths up to $n+1$,
because every raw cube degree lies between zero and $n$. Its analogous
object bound is
\begin{equation}\label{eq:exact-checkpoint-P}
 P_{\mathrm{exact}}(n,w)
 \le M(w)\left(1+N(w)\frac{n(n+3)}2\right),
\end{equation}
using $\sum_{l=2}^{n+1}l=n(n+3)/2$. Absolute birth shifts are recorded
in the exact multiplicity polynomials. The universal cutoff-three
variant for arbitrary module continuations, if desired, has the
decision bound $M(w)(1+5N(w))$ instead of~\eqref{eq:checkpoint-P}.

\subsection{Complexity through impure stretches}

\begin{theorem}[Quantified checkpoint complexity]
\label{thm:checkpoint-complexity}
Under the preceding assumptions, both the exact interval-sharing
backend and the shortened decision backend admit the conservative
bit-complexity bound
\begin{equation}\label{eq:checkpoint-runtime}
 T(n;w,g)\le n^{O(1)}
 \exp\!\left(O\!\left(g+w\log(w+1)+2^{w/2}\right)\right).
\end{equation}
The same expression is a conservative space bound. Exact splitting
operations that preserve the number of objects and cost polynomial
time in the physical matrix size and coefficient dimension may be
included without changing this form of bound.
\end{theorem}

\begin{proof}
In one crossing extension, each old loopless matching has two
smoothings. Each smoothing creates at most two new circles, since every
new circle contains one of the two newly added smoothing arcs. Delooping
therefore creates at most four objects per smoothing, or eight objects
per old object in total. Elimination removes objects; exact basis
changes and interval normalization do not increase their total number;
sharing and shortening only decrease it. Thus, between two checkpoints,
the physical object count is bounded by $8^g$ times the preceding
checkpoint bound, allowing an immaterial additional factor eight
depending on endpoint accounting.

Unit elimination uses polynomially many matrix and coefficient
operations in this physical size. The interval algorithm has
polynomial cost by~\eqref{eq:rank-invariant}--\eqref{eq:mobius-bars}.
The same is true of the exact splitting procedures admitted in the
statement. In the dot-mask algebra, morphisms have $2^{O(w)}$ coefficient
bits, and direct enumeration of monomials and their gluings gives a
conservative $2^{O(w)}$ bit cost per composition or transfer, including
the four temporary ports of the next crossing. This is the same finite
algebra cost model used in the prior analysis~\cite{prior}.

Exact multiplicity polynomials have at most $n+1$ degree positions.
The virtual unreduced allocation grows by at most eight per crossing,
so each coefficient has $O(n)$ bits. In decision mode the weights have
constant size. Complete key comparisons, matching identifiers, and
persistent caches add polynomial factors in the stated physical size
and $n$; the bound need not rely on average-case hashing.

Finally,
\[
 \log M(w)=O(w\log(w+1)),\qquad
 \log(1+N(w))=O(2^{w/2}).
\]
Substitute either~\eqref{eq:checkpoint-P} or
\eqref{eq:exact-checkpoint-P} into the physical-size estimate, charge
the polynomial operations across $n$ crossings, and absorb the
resulting fixed powers in the exponential. This gives
\eqref{eq:checkpoint-runtime}.
\end{proof}

\begin{corollary}[A specific quasi-polynomial regime]
\label{cor:checkpoint-quasipolynomial}
Suppose a diagram family has constructible scan orders and a proved
checkpoint schedule for which
\begin{equation}\label{eq:checkpoint-qp-condition}
 g+2^{w/2}+w\log(w+1)=O(\log^2(n+2)).
\end{equation}
If constructing the order and schedule costs at most $n^{O(\log n)}$
bit operations, the resulting complete recognizer runs in
$n^{O(\log n)}$ bit operations on that family. Fixed $w$ and fixed $g$
give a polynomial bound when the order is polynomial-time constructible.
\end{corollary}

For integer $w$, the last term in~\eqref{eq:checkpoint-qp-condition}
can be absorbed into a universal multiple of $2^{w/2}$; it is displayed
to retain the actual matching-count contribution. For example,
$w\le4\log_2\log(n+2)+O(1)$ and $g=O(\log^2(n+2))$ suffice.

The previous component-sharing theorem required a bound on the maximum
post-elimination connected component size $b$, with the sufficient
condition $b^2 2^{w/2}=O(\log^2 n)$~\cite{prior}. The new result removes
that component-size requirement at pure checkpoints. In particular,
an interval can have length proportional to $n$, and a connected pure
component can contain arbitrarily many copies mixed by binary changes
of basis. Neither prevents the checkpoint bound. The exact and decision
variants achieve this in different ways: exact normalization keeps
polynomially many possible lengths, while the continuation quotient
uses only two lengths.

This remains a conditional structural theorem. No argument here shows
that every knot diagram, or every prime nonhomogeneous diagram left
after the structural front end, has frequent pure checkpoints. Taking
only the initial and final checkpoints can leave $g=n$, which yields
an exponential bound. Likewise, a large frontier remains costly because
$N(w)$ is doubly exponential in $w$. Experimental success of an exact
splitter on individual diagrams is useful evidence about these
parameters, but does not replace the missing uniform theorem.

\subsection{A sharper bound for restricted coefficient alphabets}
\label{sec:coefficient-alphabets}

The all-polynomials count $N(w)$ is deliberately conservative. It can be
replaced when the coefficients occurring at the useful checkpoints have
additional structure. The restriction need only hold at those
checkpoints; arbitrary intermediate coefficients are still charged at
their full dot-mask size.

\begin{corollary}[Checkpoint coefficient alphabets]
\label{cor:coefficient-alphabet}
Suppose that, at every pure checkpoint and for each actual matching $m$,
all nonsingleton pure components use coefficients from a set
$\Theta_{i,m}$ with
\[
 |\Theta_{i,m}|\le A(w).
\]
The sets may depend on the checkpoint $i$ and matching $m$; only their
cardinality bound must be uniform. Under the other assumptions of
Theorem~\ref{thm:checkpoint-complexity}, the decision checkpoint has at
most $M(w)(1+2A(w))$ objects, and both the exact and decision variants
admit the bound
\begin{equation}\label{eq:alphabet-runtime}
 T(n;w,g,A)\le n^{O(1)}
 \exp\!\left(O\!\left(g+w\log(w+1)+\log(1+A(w))\right)\right).
\end{equation}
\end{corollary}

\begin{proof}
For each matching, retain one singleton type and at most $A(w)$
length-two types, giving the stated object count. The exact variant
again gains only a polynomial factor in $n$ from possible lengths.
Between checkpoints the same factor $8^g$ bounds expansion. Importantly,
we still charge $2^{O(w)}$ bit operations for general coefficient
compositions and transfers; we do not assume that low-complexity
coefficients remain so between checkpoints. Its logarithmic contribution
$O(w)$ is absorbed by $O(w\log(w+1))$ for positive even $w$, and the
empty-frontier case is constant. Substitution in the preceding proof
gives~\eqref{eq:alphabet-runtime}.
\end{proof}

For example, put $p=w/2$. If every checkpoint coefficient is a nonzero
\emph{linear} combination of the $p$ dot generators, then
\[
 A(w)\le2^p-1,
 \qquad
 T\le n^{O(1)}\exp O\!\left(g+w\log(w+1)\right).
\]
More generally, if every such coefficient has dot degree at most a
fixed positive integer $d$, the squarefree monomial basis gives
\begin{equation}\label{eq:degree-alphabet}
 A(w)\le2^{S_d(p)}-1,
 \qquad
 S_d(p)=\sum_{j=1}^{\min(d,p)}\binom pj.
\end{equation}
For fixed $d$, this yields
\[
 T\le n^{O(1)}\exp O_d\!\left(g+w\log(w+1)+w^d\right).
\]
In particular, either the linear or the degree-at-most-two hypothesis,
together with $w=O(\log(n+2))$ and $g=O(\log^2(n+2))$, suffices for
the target $n^{O(\log n)}$. The full coefficient arithmetic is still
polynomial in $n$ in this width regime.

These are additional, checkable checkpoint hypotheses. They are not
asserted invariants of crossing extension, elimination, or scalar
splitting, and the measured corpus does not establish them for a
uniform class of residual knot diagrams. Their value is to exhibit a
less restrictive route than counting all $2^{2^{w/2}-1}$ augmentation
polynomials: a proof controlling the coefficients at recurring
checkpoints could permit logarithmic frontier width.
